\section{Light --- Lib --- Merge}
%%% generated by the blueprint pipeline from TCSlib.Programs.Light.Lib.Merge
\subsection*{Light --- Lib --- Merge: Introduction}

Declarations of the Lean module \texttt{TCSlib.Programs.Light.Lib.Merge}.

\subsection*{Light --- Lib --- Merge: Results}

\begin{theorem}[Left step in merging]
\lean{Light.Merge.step_left}
\label{Light.Merge.step_left}
\leanok
\difficulty{3}
Let $l$ and $r$ be lists of integers and let $i,j$ be natural numbers with $i<|l|$ and
$j<|r|$. If the $i$th entry of $l$ is at most the $j$th entry of $r$, then merging the
tails of $l$ and $r$ beginning at those indices gives the $i$th entry of $l$ followed by
the merge of the tail of $l$ after that entry with the unchanged tail of $r$ beginning
at $j$.
\end{theorem}

\begin{theorem}[Right step in merging]
\lean{Light.Merge.step_right}
\label{Light.Merge.step_right}
\leanok
\difficulty{3}
Let $l$ and $r$ be lists of integers and let $i,j$ be natural numbers with $i<|l|$ and
$j<|r|$. If the $j$th entry of $r$ is strictly less than the $i$th entry of $l$, then
merging the tails beginning at indices $i$ and $j$ gives the $j$th entry of $r$ followed
by the merge of the unchanged tail of $l$ beginning at $i$ with the tail of $r$
beginning at $j+1$.
\end{theorem}

\begin{theorem}[Left step when the right list ends]
\lean{Light.Merge.step_left_end}
\label{Light.Merge.step_left_end}
\leanok
\difficulty{3}
Let $l$ and $r$ be lists of integers and let $i,j$ be natural numbers with $i<|l|$ and
$|r|\le j$. Merging the tail of $l$ beginning at index $i$ with the tail of $r$
beginning at index $j$ gives the $i$th entry of $l$ followed by the merge of the tail of
$l$ beginning at $i+1$ with that same tail of $r$.
\end{theorem}

\begin{theorem}[Right step when the left list ends]
\lean{Light.Merge.step_right_end}
\label{Light.Merge.step_right_end}
\leanok
\difficulty{3}
Let $l$ and $r$ be lists of integers and let $i,j$ be natural numbers with $|l|\le i$
and $j<|r|$. Merging the tails beginning at indices $i$ and $j$ gives the $j$th entry of
$r$ followed by the merge of the unchanged tail of $l$ beginning at $i$ with the tail of
$r$ beginning at $j+1$.
\end{theorem}

\begin{definition}[Constant zero list index]
\lean{Light.Merge.ListA}
\label{Light.Merge.ListA}
\leanok
The abbreviation $	ext{ListA}$ denotes the natural number $0$.
\end{definition}

\begin{definition}[Value of the second length constant]
\lean{Light.Merge.LenA}
\label{Light.Merge.LenA}
\leanok
\uses{Light.Merge.ListA}
Define the natural number $\mathrm{LenA}$ to be $1$.
\end{definition}

\begin{definition}[Second list base address]
\lean{Light.Merge.ListB}
\label{Light.Merge.ListB}
\leanok
\uses{Light.Merge.ListA}
The base address of the second list is the natural number $2$.
\end{definition}

\begin{definition}[Second list length address]
\lean{Light.Merge.LenB}
\label{Light.Merge.LenB}
\leanok
\uses{Light.Merge.ListA}
The address assigned to the length of the second list is the natural number $3$.
\end{definition}

\begin{definition}[Destination address]
\lean{Light.Merge.Dest}
\label{Light.Merge.Dest}
\leanok
\uses{Light.Merge.ListA}
The destination address is the natural number $4$.
\end{definition}

\begin{definition}[First completion counter address]
\lean{Light.Merge.DoneA}
\label{Light.Merge.DoneA}
\leanok
\uses{Light.Merge.ListA}
The address assigned to the completion counter for the first list is the natural number
$5$.
\end{definition}

\begin{definition}[Second completion counter address]
\lean{Light.Merge.DoneB}
\label{Light.Merge.DoneB}
\leanok
\uses{Light.Merge.ListA}
The address assigned to the completion counter for the second list is the natural number
$6$.
\end{definition}

\begin{definition}[Left step of a merge]
\lean{Light.mergeTakeLeft}
\label{Light.mergeTakeLeft}
\leanok
\uses{Light.M, Light.Merge.Dest, Light.Merge.DoneA, Light.Merge.DoneB, Light.Merge.ListA, Light.Stmt, Light.k, Light.v}
The statement first stores the value read from address $	ext{ListA}+	ext{DoneA}$ at
address $	ext{Dest}+	ext{DoneA}+	ext{DoneB}$. It then increases $	ext{DoneA}$ by one.
\end{definition}

\begin{definition}[Right-List Merge Step]
\lean{Light.mergeTakeRight}
\label{Light.mergeTakeRight}
\leanok
\uses{Light.M, Light.Merge.Dest, Light.Merge.DoneA, Light.Merge.DoneB, Light.Merge.ListB, Light.Stmt, Light.k, Light.v}
The program stores the value of the second list at the position obtained by adding the
destination base address and the two completion counters. It then increases the second
completion counter by one.
\end{definition}

\begin{definition}[Merge loop body]
\lean{Light.mergeBody}
\label{Light.mergeBody}
\leanok
\uses{Light.M, Light.Merge.DoneA, Light.Merge.DoneB, Light.Merge.LenA, Light.Merge.LenB, Light.Merge.ListA, Light.Merge.ListB, Light.Stmt, Light.k, Light.mergeTakeLeft, Light.mergeTakeRight, Light.v}
The body sets both completion counters to zero and then repeats while their sum is less
than the sum of the two list lengths. When both lists have remaining elements, it takes
the next element from the second list if that element is smaller than the next element
from the first list, and otherwise takes the next element from the first list; when only
one list has remaining elements, it takes the next element from that list.
\end{definition}

\begin{definition}[Merge preconditions]
\lean{Light.Merge.Pre}
\label{Light.Merge.Pre}
\leanok
\uses{Light.Apart, Light.Limits, Light.Seg}
For lists $L,R$ of integers, memory $\mu$, and addresses $a,b,d$, the merge precondition
requires that memory stores $L$ and $R$ on the segments beginning at $a$ and $b$, all
input and destination segments fit within the space bound, the space bound is no greater
than the word bound, and each input segment is disjoint from the destination segment of
length $|L|+|R|$ beginning at $d$.
\end{definition}

\begin{definition}[Merge loop invariant]
\lean{Light.Merge.Inv}
\label{Light.Merge.Inv}
\leanok
\uses{Light.SameOutside, Light.Seg, Light.State, Light.frame}
For lists $L,R$, memory $\mu$, addresses $a,b,d$, and indices $i,j$, the merge invariant
says that the state has local frame $[a,|L|,b,|R|,d,i,j]$ and some memory $\mu\prime$;
$i\le |L|$ and $j\le |R|$; $\mu\prime$ agrees with $\mu$ outside the destination
segment of length $|L|+|R|$; and an output list of length $i+j$ is stored at $d$ such
that this output followed by the merge of the unconsumed suffixes $L.drop(i)$ and
$R.drop(j)$ equals $L.merge(R)$.
\end{definition}

\begin{definition}[Number of elements already taken]
\lean{Light.Merge.Taken}
\label{Light.Merge.Taken}
\leanok
\uses{Light.Merge.Inv, Light.State}
For lists $L,R$, initial memory $\mu$, addresses $a,b,d$, a natural number $t$, and state $\sigma$,
the state has taken $t$ elements when there exist natural numbers $i,j$ with $i+j=t$ such
that $\sigma$ satisfies the merge invariant at indices $i,j$.
\end{definition}

\begin{theorem}[Merge step taking the left entry]
\lean{Light.Merge.takeLeft}
\label{Light.Merge.takeLeft}
\leanok
\uses{Light.Merge.Inv, Light.Merge.Pre, Light.Merge.Taken, Light.SameOutside.update, Light.Seg.get, Light.Seg.snoc, Light.State, Light.Stmt.Runs, Light.mergeTakeLeft, Light.update_frame_setLocal}
\difficulty{5}
Let $L$ and $R$ be integer lists, let $\mu$ be an initial memory map, and let $a,b,d$ be
natural-number source and destination addresses. Suppose the merge preconditions hold:
the space bound is at most the word bound, each list is represented in memory at its
source address, the source and destination regions fit within the space bound, and each
source region is apart from the destination region. If a state satisfies the merge
invariant at indices $i,j$, where $i<|L|$, and the merge of the remaining tails begins
with $L_i$ followed by the merge after advancing the left index, then the left-taking
merge step is block-safe and its resulting state satisfies the merge invariant for some
indices whose sum is $i+j+1$.
\end{theorem}

\begin{theorem}[Merge step taking the right entry]
\lean{Light.Merge.takeRight}
\label{Light.Merge.takeRight}
\leanok
\uses{Light.Merge.Inv, Light.Merge.Pre, Light.Merge.Taken, Light.SameOutside.update, Light.Seg.get, Light.Seg.snoc, Light.State, Light.Stmt.Runs, Light.mergeTakeRight, Light.update_frame_setLocal}
\difficulty{5}
Let $L$ and $R$ be integer lists, let $\mu$ be an initial memory map, and let $a,b,d$ be
natural-number source and destination addresses. Suppose the merge preconditions hold:
the space bound is at most the word bound, each list is represented in memory at its
source address, the source and destination regions fit within the space bound, and each
source region is apart from the destination region. If a state satisfies the merge
invariant at indices $i,j$, where $j<|R|$, and the merge of the remaining tails begins
with $R_j$ followed by the merge after advancing the right index, then the right-taking
merge step is block-safe and its resulting state satisfies the merge invariant for some
indices whose sum is $i+j+1$.
\end{theorem}

\begin{definition}[Merge running-time bound]
\lean{Light.mergeTime}
\label{Light.mergeTime}
\leanok
For every natural number $n$, define the merge time to be $40n+12$.
\end{definition}

\begin{theorem}[Correctness of the merge procedure]
\lean{Light.merge_meets}
\label{Light.merge_meets}
\leanok
\uses{Light.Ends.whileBlock, Light.Meets, Light.Meets.of_body, Light.Merge.Pre, Light.Merge.Taken, Light.Merge.step_left, Light.Merge.step_left_end, Light.Merge.step_right, Light.Merge.step_right_end, Light.Merge.takeLeft, Light.Merge.takeRight, Light.SameOutside, Light.SameOutside.refl, Light.Seg, Light.Seg.get, Light.Seg.nil, Light.Stmt.Runs.ite_neg, Light.Stmt.Runs.ite_pos, Light.mergeBody, Light.mergeTakeLeft, Light.mergeTakeRight, Light.mergeTime}
\difficulty{5}
Let a program $P$ contain the merge procedure at position $p$, and let the machine
limits have word bound $W$, memory size $S$, and depth bound $D$. Suppose the input
memory $\mu$ stores lists $l,r$ of integers starting at addresses $a,b$, respectively;
both lists and a destination segment of length $|l|+|r|$ fit within memory; the
destination segment is disjoint from each input segment; and $S\le W$. Then, when the
procedure is called at depth $d$ with inputs $(a,|l|,b,|r|,\mathrm{dst})$, it can
execute within time $40(|l|+|r|)+12$ and leave a memory $\mu'$ in which the segment
starting at $\mathrm{dst}$ contains the sorted merge of $l$ and $r$, while every memory
location outside that destination segment retains its initial value.
\end{theorem}
